\documentclass[10pt]{article}
\usepackage[margin=0.9in]{geometry}
\usepackage{amsmath,amssymb}
\usepackage[T1]{fontenc}
\usepackage[utf8]{inputenc}
\usepackage{lmodern}
\usepackage{microtype}
\usepackage{booktabs}
\usepackage{enumitem}
\usepackage[hidelinks]{hyperref}
\hypersetup{pdftitle={A Two-Configuration Cavendish Test That Separates Low-Acceleration Gravity Models: A Proposal for G Metrology Groups}}
\setlist{nosep,leftmargin=1.3em}
\setlength{\parskip}{4pt}\setlength{\parindent}{0pt}
\sloppy
\begin{document}
\begin{center}
{\Large\bfseries A Two-Configuration Cavendish Test That Separates\\ Low-Acceleration Gravity Models}\\[4pt]
{\large A proposal for G metrology groups}\\[6pt]
{\small Working note, 2026-10-15, v1.0. AI-assisted research programme; not peer reviewed.}
\end{center}

\begin{abstract}\noindent
Low-acceleration modifications of gravity make sharply different predictions for one inexpensive laboratory measurement. That measurement is the fractional change in the measured $G$ when the gravitational acceleration that the source masses produce at the test masses is changed by a factor of four, with everything else held fixed.

\begin{itemize}
\item \textbf{Source-field-boost model.} One published model applies a MOND-type boost to the source field alone (Klein's interpolation, $\beta\simeq1.3$, as used by Scarpa, Falomo \& Treves 2026). It predicts a shift of about 340\,ppm between configurations at half and twice the source acceleration of the Quinn et al.\ apparatus.
\item \textbf{Dark-energy-tied model.} A framework in which the low-acceleration scale is tied to dark energy and screened at laboratory scales predicts no measurable shift: $|\Delta G/G|<10^{-6}$, with a model value of order $10^{-20}$.
\item \textbf{Standard physics.} Newtonian gravity and general relativity predict zero.
\end{itemize}

Distinguishing the first prediction from the others needs a precision on the \emph{difference} of about 30--50\,ppm. Leading torsion-balance groups reach this today. We describe how to run the test on an existing apparatus: the two configurations, the systematics budget, blinding, and null checks. We also give the predictions in a form that can be checked directly.
\end{abstract}

\section{The question and the competing predictions}
\begin{itemize}
\item \textbf{The setup.} In a Cavendish-type experiment, the source masses produce an acceleration $a_s$ at the test masses. Typically $a_s\sim10^{-8}$--$10^{-7}$\,m\,s$^{-2}$, far below Earth's $g\approx9.8$\,m\,s$^{-2}$ but still well above the galactic scale $a_0\sim10^{-10}$\,m\,s$^{-2}$.
\item \textbf{The observable.} Let $G(a_s)$ be the value an experiment infers.
\item \textbf{The proposed measurement.} Measure
$$\delta\equiv\frac{G(a_1)-G(a_2)}{\bar G},\qquad a_2/a_1\approx4,$$
on one instrument, with only the source configuration changed.
\end{itemize}

\begin{center}
\begin{tabular}{lcl}
\toprule
model & predicted $\delta$ (half $\to$ twice the nominal $a_s$) & reason\\
\midrule
Newton / general relativity & 0 & $G$ is a constant\\
Source-field boost (Klein, $\beta=1.3$) [1] & $\approx3.4\times10^{-4}$ & boost $F(a_s/a_0)$ applied to the source field\\
Simple-$\nu$ boost on the source field alone (illustrative) & $\sim3\times10^{-3}$ & as above, steeper function\\
Dark-energy-tied, unscreened, Earth's field included [2] & $\sim4\times10^{-20}$ & Earth's $g$ sets the regime\\
Dark-energy-tied, committed (stiffened) form [2,3] & $0$ (offset $\sim10^{-36}$) & screened below parsec scales\\
\bottomrule
\end{tabular}
\end{center}

\textbf{Reproducing the source-boost numbers.} Klein's function is $F(x)=[1+x^{-\beta}]^{1/(2\beta)}$, with $x=a_s/a_0$, $a_0=1.2\times10^{-10}$\,m\,s$^{-2}$ and $\beta=1.3$. It reproduces the published predictions $G=6.6769$ and $6.6747\times10^{-11}$ (in SI units) for half and twice the Quinn source acceleration. That corresponds to $a_s\approx4.7\times10^{-8}$\,m\,s$^{-2}$ nominal, which is consistent with four 11\,kg source masses at about 0.21\,m. The difference is 338\,ppm. (Lane CFG772 [2] reproduces both numbers to $8\times10^{-6}$.)

\textbf{Why the dark-energy-tied framework predicts nothing.} Three independent reasons, all in [2]:
\begin{itemize}
\item Earth's $g\gg a_0$ puts the whole laboratory deep in the Newtonian regime.
\item The source masses are supported, not free-falling, so they do not define their own low-acceleration region.
\item The committed form of the theory screens its extra field on scales below about 2\,pc. This is required by the 2026 Cassini bound [4].
\end{itemize}
A source-dependent shift $|\delta|\ge10^{-6}$ would falsify that framework.

\section{How to run it on an existing apparatus}
Any of the three standard readouts works: time-of-swing, angular acceleration feedback, or electrostatic servo. The proposal adds no new hardware beyond a second source configuration.

\subsection*{Choosing the two configurations}
\begin{itemize}
\item \textbf{Option A, change the source mass (preferred).} Replace the source masses with a geometrically identical set of four times the density, or four times the mass at the same centre of mass. Examples: tungsten-alloy versus copper inserts, or hollow versus solid shells. Positions are unchanged, so the dominant geometric terms in the $G$ extraction are unchanged. What is new is the density metrology (mass, homogeneity, dimensions) of the second set.
\item \textbf{Option B, change the source distance.} Move the same masses so that $a_s$ at the test masses changes by 4. This needs metrology of the new positions and a recomputed geometry model. Cross-check the multipole coefficients against a dedicated calibration run.
\item \textbf{Option C, both.} Run A and B together and require the same $G$ from all three to four configurations. Combining them separates any source-acceleration dependence from geometry errors.
\end{itemize}

\subsection*{Measurement plan}
\begin{enumerate}
\item \textbf{Interleaving.} Alternate the two configurations in blocks (e.g.\ ABBA over days to weeks). This cancels slow drifts in fibre properties, temperature and tilt.
\item \textbf{Blinding.} An external party adds a hidden offset to the inferred $G$ of each configuration. The offsets are removed only after the systematics budget is frozen.
\item \textbf{Null checks.} Rotate the whole source assembly and reverse its orientation. Use a ``dummy'' configuration with the same $a_s$ but different masses: it must give $\delta=0$. Vary the vacuum pressure and the temperature setpoint.
\item \textbf{Statistic.} Report $\delta$ with a full covariance, including correlated systematics between configurations.
\end{enumerate}

\subsection*{Precision budget}
\begin{itemize}
\item \textbf{Target.} $\sigma(\delta)\le50$\,ppm separates the source-boost prediction from zero at $\ge6\sigma$. $\sigma(\delta)\le30$\,ppm gives $\ge10\sigma$.
\item \textbf{Why this is achievable.} The difference measurement inherits only the systematics that \emph{change} between configurations: mass and density metrology (Option A) or position metrology (Option B). Absolute $G$ measurements already reach about 10--30\,ppm per apparatus [5], so a differential $\sigma(\delta)\sim20$--40\,ppm is a realistic first goal.
\item \textbf{Reaching $10^{-6}$.} The full kill line for the dark-energy-tied framework needs about 20--30$\times$ beyond today's best. That would be a longer-term programme.
\end{itemize}

\section{What each outcome would mean}
\begin{itemize}
\item \textbf{$|\delta|\approx3\times10^{-4}$ (source-dependent).} Laboratory $G$ depends on the source acceleration. That excludes Newtonian gravity at laboratory scales, excludes the dark-energy-tied framework, and supports source-field-boost models.
\item \textbf{$|\delta|<50$\,ppm.} Source-field-boost models of the published size are excluded. The null predictions survive. The discrepancy between published laboratory $G$ values then has to come from apparatus systematics.
\item \textbf{Any intermediate value.} It would be the first measured acceleration dependence of laboratory $G$, and a strong result whatever the theory.
\end{itemize}
The outcome is informative either way, at modest cost: one additional source configuration on an apparatus that already exists.

\section{Caveats}
\begin{itemize}
\item The Quinn geometry used here (11\,kg sources at about 0.21\,m) and the Scarpa et al.\ numbers were transcribed from memory and from a summary of their paper. They are approximate; each group should recompute the predictions for its own geometry.
\item The dark-energy-tied predictions come from a framework with a fitted coefficient ($\kappa=\tfrac12$) and a bounded stiffness length. They are falsifiable predictions, not established physics.
\end{itemize}

\vspace{4pt}\hrule\vspace{3pt}
{\footnotesize
[1] R.~Scarpa, R.~Falomo \& A.~Treves, ``Gravity anomaly from laboratory experiments to astrophysics'', arXiv:2603.17770 (2026), using Klein's interpolation.
[2] This programme, lane CFG772 (frozen criteria, script and outputs in the public repository), with the sealed kill line $|\Delta G/G|\ge10^{-6}$.
[3] Lane CFG770: the stiffened action against the 2026 Cassini quadrupole bound.
[4] R.~S.~Park et al., arXiv:2602.17884 (2026).
[5] T.~Quinn et al., Phys.\ Rev.\ Lett.\ 111, 101102 (2013); CODATA 2018 G compilation.
AI assistance: Anthropic's Claude.}
\end{document}
